\chapter{THIẾT KẾ LOGIC}

\section{Ánh xạ ERD sang lược đồ quan hệ}

Áp dụng quy tắc ánh xạ chuẩn: thực thể mạnh → một quan hệ; mối kết hợp 1:N → khóa chính phía 1 đưa sang phía N làm khóa ngoại; mối kết hợp M:N → tạo quan hệ trung gian với khóa chính là tổ hợp hai khóa ngoại; thực thể yếu → khóa chính là khóa chủ ghép khóa cục bộ.

\begin{infobox}
\footnotesize
\begin{verbatim}
NguoiDung(maND, email, matKhauHash, maTV FK, ngayTao)
GiaDinh(maGD, tenGiaDinh, ngayTao)
ThanhVien(maTV, maGD FK, hoTen, ngaySinh, ngayMat, gioiTinh)
QuanHeChaCon(maCon FK, maChaMe FK, vaiTro) -- PK: (maCon, maChaMe)
HonNhan(maHN, maTV1 FK, maTV2 FK, ngayKetHon, ngayLyHon, trangThai)
VaiTro(maVT, tenVT)
PhanQuyen(maND FK, maGD FK, maVT FK) -- PK tổ hợp 3 khóa
BaiViet(maBV, maND FK, maGD FK, noiDung, ngayDang)
BinhLuan(maBV FK, sttBL, maND FK, noiDung, ngayBL) -- khóa: (maBV, sttBL)
SuKien(maSK, maGD FK, tenSK, diaDiem, ngayDienRa)
ThamGiaSuKien(maSK FK, maND FK, trangThai, ngayThamGia) -- PK: (maSK, maND)
HoSoNgheNghiep(maHS, maTV FK, ngheNghiep, noiCongTac, tuNam, denNam)
TaiLieuLichSu(maTL, maGD FK, tieuDe, loaiTepTin, url)
QuanTriGiaDinh(maND PK/FK, quyenDacBiet) -- ánh xạ chuyên biệt hóa 3.4,
                                            1:1 với NguoiDung
ThanhVienThuong(maND PK/FK) -- ánh xạ chuyên biệt hóa 3.4
NhatKyKiemToan(maNK, maND FK NULL, hanhDong, doiTuong, thoiGian)
\end{verbatim}
\end{infobox}

\section{Xác định phụ thuộc hàm (FD)}

Phụ thuộc hàm được rút ra trực tiếp từ quy tắc nghiệp vụ ở mục 2.2, làm cơ sở để kiểm tra và chứng minh chuẩn hóa:

\begin{infobox}
\begin{verbatim}
maTV -> maGD, hoTen, ngaySinh, ngayMat, gioiTinh
maND -> email, matKhauHash, maTV
(maCon, maChaMe) -> vaiTro
(maSK, maND) -> trangThai, ngayThamGia
maHS -> maTV, ngheNghiep, noiCongTac, tuNam, denNam
maBV -> maND, maGD, noiDung, ngayDang
(maBV, sttBL) -> maND, noiDung, ngayBL
maHN -> maTV1, maTV2, ngayKetHon, ngayLyHon, trangThai
maVT -> tenVT
maTL -> maGD, tieuDe, loaiTepTin, url
maNK -> maND, hanhDong, doiTuong, thoiGian
\end{verbatim}
\end{infobox}

\section{Kiểm tra và chứng minh chuẩn hóa}

\subsection{Dạng chuẩn 1 (1NF)}

Không có thuộc tính đa trị được lưu chung trong một cột. Ví dụ, số điện thoại liên hệ của ThanhVien nếu có nhiều giá trị sẽ được tách thành bảng riêng:

\begin{infobox}
\begin{verbatim}
LienHeThanhVien(maTV FK, loaiLienHe, giaTri)
  -- PK: (maTV, loaiLienHe, giaTri)
\end{verbatim}
\end{infobox}

\subsection{Dạng chuẩn 2 (2NF)}

Xét các quan hệ có khóa tổ hợp: ThamGiaSuKien(maSK, maND) → trangThai, ngayThamGia; QuanHeChaCon(maCon, maChaMe) → vaiTro; PhanQuyen(maND, maGD, maVT). Mọi thuộc tính không khóa đều phụ thuộc đầy đủ vào toàn bộ khóa tổ hợp, không có thuộc tính nào chỉ phụ thuộc một phần khóa → các quan hệ đạt 2NF.

Các quan hệ đơn khóa bổ sung ở 4.2 (HonNhan, VaiTro, TaiLieuLichSu, NhatKyKiemToan, QuanTriGiaDinh, ThanhVienThuong) đạt 2NF hiển nhiên vì khóa chính không phải tổ hợp. Riêng PhanQuyen không có thuộc tính không khóa nào nên điều kiện 2NF cũng thỏa mãn một cách hiển nhiên (không có gì để kiểm tra phụ thuộc bộ phận).

\subsection{Dạng chuẩn 3 (3NF) --- chứng minh bằng phản ví dụ}

Giả sử ban đầu gộp ThanhVien và GiaDinh thành một bảng ThanhVien\_GD(maTV, hoTen, maGD, tenGiaDinh). Khi đó tồn tại phụ thuộc bắc cầu:

\begin{center}
maTV → maGD, \quad maGD → tenGiaDinh \quad $\Rightarrow$ \quad maTV → tenGiaDinh \ (bắc cầu)
\end{center}

Điều này gây dị thường: đổi tên một GiaDinh phải sửa trên mọi dòng ThanhVien thuộc gia đình đó (bất thường sửa), và không thể tạo một GiaDinh mới nếu chưa có ThanhVien nào (bất thường thêm). Phân rã tách GiaDinh thành quan hệ riêng như ở mục 4.1 loại bỏ phụ thuộc bắc cầu này → đạt 3NF.

\subsection{Dạng chuẩn Boyce--Codd (BCNF)}

Với mọi FD không tầm thường X → Y trong các quan hệ đã liệt kê, vế trái X đều là siêu khóa (khóa chính hoặc khóa chính tổ hợp) của quan hệ tương ứng. Do đó toàn bộ lược đồ ở mục 4.1 đạt BCNF; không phát sinh phụ thuộc bất thường ngoài khóa.

Kết luận này nay bao quát toàn bộ các quan hệ đã liệt kê ở 4.1, kể cả các bảng bổ sung HonNhan, VaiTro, TaiLieuLichSu, NhatKyKiemToan, QuanTriGiaDinh và ThanhVienThuong --- tất cả đều có khóa chính là khóa đơn hoặc khóa ngoại 1:1, nên mọi FD không tầm thường đều xuất phát từ siêu khóa.

\section{Kiểm tra tính bảo toàn của phép phân rã}

Áp dụng điều kiện đủ (R1 $\cap$ R2) → R1 hoặc (R1 $\cap$ R2) → R2: ví dụ ThanhVien và GiaDinh giao nhau ở maGD, là khóa chính của GiaDinh → phép phân rã bảo toàn kết nối. Tương tự, BaiViet và BinhLuan giao nhau ở maBV, là khóa chính của BaiViet → bảo toàn kết nối, không sinh dòng giả khi JOIN lại.

